% CLAUDE REPORT::WCDERIV-0001::2026-11-06
\documentclass{article}
\usepackage{amsmath,amssymb,amsthm}
\begin{document}
\title{Worst-case derivation complexity $\overline W_R(N)$: bridge audit (WCDERIV-0001)}
\maketitle
\noindent\textbf{Status: P versus NP remains OPEN.} No lower bound or bridge theorem is proved here.

\section*{1. Reduction of the directive to known proof complexity}
Free choice of the optimal derivation is nondeterminism, so $\overline W_R(N)$ is the Cook--Reckhow proof-size function of $R$ (with the cost measure replacing size).
\begin{itemize}
\item \textbf{KNOWN THEOREM (Cook--Reckhow 1979).} A polynomially bounded sound propositional proof system exists iff $\mathrm{NP}=\mathrm{coNP}$.
\item \textbf{PROVED/EXACT (consequence).} For a fixed sound $R$, $\overline W_R$ polynomial $\Rightarrow$ $\mathrm{NP}=\mathrm{coNP}$. A superpolynomial bound for one fixed $R$ implies nothing about $\mathrm{P}\ne\mathrm{NP}$ by itself; the bridge is therefore load-bearing.
\item \textbf{PROVED/EXACT.} If $\mathrm{SAT}\in\mathrm{P}$ then $\mathrm{NP}=\mathrm{coNP}$, hence some $R^\ast$ has $\overline W_{R^\ast}\le \mathrm{poly}$. So the bridge for \emph{some} system is free; the bridge for a \emph{fixed, named} $R$ is the content.
\end{itemize}

\section*{2. Exact form of the bridge for a fixed $R$}
\textbf{Proposition (PROVED/EXACT, short).} Let $\mathrm{Bridge}(R)$ denote $\mathrm{SAT}\in\mathrm P\Rightarrow\overline W_R\le\mathrm{poly}$.
(a) If $R$ is $p$-optimal or merely optimal (simulates every system), $\mathrm{Bridge}(R)$ holds: apply the first bullet's $R^\ast$.
(b) $\mathrm{Bridge}(R)$ together with $\overline W_R\ne\mathrm{poly}$ gives $\mathrm P\ne\mathrm{NP}$ (contrapositive).
(c) Conversely any proof of $\mathrm{Bridge}(R)$ for a system not known optimal must use that $R$ can certify, for the solver $M$ supplied by the hypothesis, the step ``$M$ rejects $F$ $\Rightarrow$ $F$ unsat''. A poly-size transcript of $M$ is a valid object, but its soundness is a statement about all assignments, which $R$ has to prove. \textbf{CANDIDATE:} for $R=$ Extended Frege, if PV (or $S^1_2$) proves correctness of $M$, EF has poly-size proofs (Cook's translation). This is a \emph{provability} hypothesis, strictly stronger than $\mathrm{SAT}\in\mathrm P$, so it does not give the unconditional bridge.

\section*{3. Consequences for the five roles}
\begin{enumerate}
\item \textbf{Lower-bound lead:} the weakest sufficient bridge is $\mathrm{Bridge}(R)$ for an $R$ in which hard instances are provable to be hard; for $R$ weaker than EF (resolution, cutting planes, bounded-depth Frege) $\mathrm{Bridge}(R)$ is \emph{false-looking}: resolution is automatable only in quasi-poly size terms and exponential lower bounds exist unconditionally, so $\mathrm{Bridge}(\mathrm{Res})$ is equivalent to $\mathrm{P}\ne\mathrm{NP}$ being false \,---\, i.e.\ it cannot be a theorem. \textbf{PROVED/EXACT:} known exponential resolution lower bounds (Haken) refute $\mathrm{Bridge}(\mathrm{Res})$ conditional on $\mathrm{P=NP}$ only if such lower bounds are unconditional, which they are; hence $\mathrm{Bridge}(\mathrm{Res})$ is equivalent to $\mathrm{P}\ne\mathrm{NP}$, a vacuous bridge. The same applies to every system with an unconditional superpoly lower bound. The useful $R$ must have \emph{no} known lower bound: Frege, EF, and stronger.
\item \textbf{Circuit-reduction builder:} a fixed combined inference system strong enough for $\mathrm{Bridge}(R)$ must simulate EF, so shortest-DAG certification must target EF-style (extension-variable) derivations; lengths are then Frege/EF lower-bound questions, open.
\item \textbf{Hard-UNSAT family:} candidates must survive EF. Known: no superpolynomial EF lower bound is known for any family; Pich--Santhanam give evidence (conditional on hardness assumptions) that natural random-formula distributions cannot be proved hard for strong systems.
\item \textbf{Resource auditor:} space/time cost may only increase $\overline W$; any cost model weaker than size is covered by the same argument.
\item \textbf{Red team:} the dichotomy above is the sharp counter-check: any claimed bridge for a system with a known unconditional lower bound is vacuous or wrong.
\end{enumerate}

\section*{4. Literature check (this run)}
Pich--Santhanam, FOCS 2019 (barriers to proof complexity lower bounds for strong systems, conditional on cryptographic/Rudich-type assumptions): \textbf{useful structural analogy}; it is evidence the EF lower bound is hard, not a proof. Cook--Reckhow and the optimal-system literature: \textbf{directly transferable theorem}, used in \S1--2. Optimal proof systems exist under $\mathrm{NE}=\mathrm{coNE}$; EF with soundness axioms simulates any system; unconditional existence of an optimal system is open.

\section*{5. Corpus audit}
\textbf{CORPUS REVIEWED:} AGENTS.md, repository tree, research/agents/claude (world ledger, TPORT, ONETAPE), a repository grep for Cook--Reckhow/worst-case derivation (LSPV-1/350, LSPV-5/163 proof-system fallback, grok PROOFSPACE1). Not exhaustively read; MAIN snapshots 2026-11-05 only listed. \textbf{DUPLICATION:} LSPV-5/163 and LSPV-1/350 touch proof-system fallbacks and worst-case quotients; the dichotomy in \S3.1 was not checked against them line by line. \textbf{WHAT CHANGED RELATIVE TO ROOT MAIN:} nothing; root \texttt{main.tex} was not edited (this role owns no region). \textbf{STALE CLAIMS FOUND:} none checked. \textbf{ADOPTED/REJECTED:} none.

\section*{6. Next step}
Decide a single target system between Frege and EF and check whether any corpus hard family (LSPV-1/2 signed hypergraph instances) is even a candidate against EF; otherwise the project's lower-bound target is vacuous or open. Status: bridge for strong $R$ OPEN; lower bound OPEN.
\end{document}
% END_OF_VERIFIED_REPORT::CLAUDE::WCDERIV-0001
